% =====================================================================
%  BÁO CÁO ĐỒ ÁN MÔN HỌC -- CSE703153 (Software Security)
%  Đề tài: OrderPaymentMicroservice -- Nhóm 9
%
%  BIÊN DỊCH: khuyến nghị XeLaTeX hoặc LuaLaTeX (DejaVu Serif, tiếng Việt đầy đủ).
%  Khối iftex bên dưới tự chọn cấu hình pdfLaTeX (T5) hoặc Xe/LuaLaTeX (fontspec);
%  pdfLaTeX cần bộ font T5/vntex đầy đủ.
%
%  HÌNH ẢNH: đặt các tệp anh-01..anh-10 (.png) cùng thư mục với file .tex.
%  Macro \anh{ten-tep}{caption} tự chèn ảnh khi tệp tồn tại; nếu chưa có,
%  in khung nhắc để báo cáo vẫn biên dịch được.
%
%  ĐIỀN TRƯỚC KHI NỘP: mã lớp; bổ sung thành viên nếu nhóm có nhiều người.
% =====================================================================
\documentclass[12pt,a4paper]{article}

% Khuyến nghị XeLaTeX/LuaLaTeX (DejaVu Serif, hỗ trợ tiếng Việt đầy đủ).
% Nếu buộc dùng pdfLaTeX, cần bộ font T5/vntex; khối iftex dưới đây sẽ chọn
% cấu hình tương ứng.
\usepackage{iftex}
\ifPDFTeX
  \usepackage[utf8]{inputenc}
  \usepackage[T5]{fontenc}
  \usepackage[vietnamese]{babel}
\else
  \usepackage{fontspec}
  % Nếu máy không có DejaVu, đổi sang Times New Roman hoặc Latin Modern Roman.
  \setmainfont{DejaVu Serif}[Scale=0.98]
  \setsansfont{DejaVu Sans}[Scale=0.98]
  \setmonofont{DejaVu Sans Mono}[Scale=0.86]
  \usepackage[vietnamese]{babel}
\fi

\usepackage[margin=2cm,headheight=15pt]{geometry}
\usepackage{titlesec}
\usepackage{titletoc}
\usepackage{listings}
\usepackage{xcolor}
\usepackage{fancyhdr}
\usepackage{url}
\usepackage{hyperref}
\usepackage{amsmath,amssymb}
\usepackage{enumitem}
\usepackage{tabularx}
\usepackage{booktabs}
\usepackage{longtable}
\usepackage{graphicx}
\usepackage{float}
\usepackage{array}
\newcolumntype{L}[1]{>{\raggedright\arraybackslash}p{#1}}
\newcolumntype{Y}{>{\raggedright\arraybackslash}X}
\usepackage{caption}
\usepackage{tikz}
\usetikzlibrary{arrows.meta,positioning,fit,calc}

% ---------- Màu và định dạng ----------
\definecolor{primary}{RGB}{0,51,102}
\definecolor{secondary}{RGB}{204,85,0}
\definecolor{codebg}{RGB}{248,249,250}
\definecolor{codegray}{RGB}{108,117,125}
\definecolor{codecomment}{RGB}{40,167,69}
\definecolor{codekeyword}{RGB}{0,86,179}
\definecolor{codestring}{RGB}{214,51,132}
\definecolor{lightline}{RGB}{220,224,230}

\pagestyle{fancy}
\fancyhf{}
\fancyhead[L]{\small\color{codegray}OrderPaymentMicroservice -- CSE703153}
\fancyhead[R]{\small\color{codegray}Nhóm 9 -- Phenikaa University}
\fancyfoot[C]{\thepage}
\renewcommand{\headrulewidth}{0.5pt}
\renewcommand{\footrulewidth}{0pt}

\titleformat{\section}
  {\normalfont\Large\bfseries\color{primary}}{\thesection}{1em}{}
\titleformat{\subsection}
  {\normalfont\large\bfseries\color{primary!90!black}}{\thesubsection}{1em}{}
\titleformat{\subsubsection}
  {\normalfont\normalsize\bfseries\color{secondary}}{\thesubsubsection}{1em}{}

\setlist[itemize,1]{label=--,leftmargin=1.4em}
\setlist[enumerate,1]{leftmargin=1.6em}

\lstset{
    backgroundcolor=\color{codebg},
    basicstyle=\ttfamily\footnotesize,
    keywordstyle=\color{codekeyword}\bfseries,
    commentstyle=\color{codecomment}\itshape,
    stringstyle=\color{codestring},
    numbers=left,
    numberstyle=\tiny\color{codegray},
    stepnumber=1,
    numbersep=8pt,
    frame=single,
    rulecolor=\color{lightline},
    breaklines=true,
    breakatwhitespace=true,
    tabsize=4,
    showspaces=false,
    showstringspaces=false,
    captionpos=b
}
\lstdefinelanguage{json}{
    basicstyle=\ttfamily\footnotesize,
    numbers=left,
    numberstyle=\tiny\color{codegray},
    showstringspaces=false,
    breaklines=true,
    frame=single,
    backgroundcolor=\color{codebg},
    rulecolor=\color{lightline},
    morestring=[b]",
    stringstyle=\color{codestring}
}
\lstdefinelanguage{bash}{
    morekeywords={python3,dotnet,make,gcc,cbmc,z3,valgrind},
    sensitive=true,
    morecomment=[l]{\#},
    morestring=[b]"
}
\lstdefinelanguage{CSharp}{
    keywords={abstract,as,async,await,base,bool,break,byte,case,catch,char,checked,class,const,continue,decimal,default,delegate,do,double,else,enum,event,explicit,extern,false,finally,fixed,float,for,foreach,goto,if,implicit,in,int,interface,internal,is,lock,long,namespace,new,null,object,operator,out,override,params,private,protected,public,readonly,record,ref,return,sbyte,sealed,short,sizeof,stackalloc,static,string,struct,switch,this,throw,true,try,typeof,uint,ulong,unchecked,unsafe,ushort,using,var,virtual,void,volatile,when,while,nameof},
    sensitive=true,
    morecomment=[l]{//},
    morecomment=[s]{/*}{*/},
    morestring=[b]"
}

\hypersetup{
    colorlinks=true,
    linkcolor=primary,
    urlcolor=primary,
    citecolor=primary,
    hypertexnames=false
}
\emergencystretch=3em

% ---------- Macro tiện dụng ----------
\newcommand{\cmd}[1]{\texttt{\detokenize{#1}}}
\newcommand{\cwe}[1]{CWE-#1}
\newcommand{\mustfill}[1]{\textcolor{secondary}{\textbf{[#1]}}}

% Ảnh minh chứng. Cú pháp: \anh[độ rộng]{tên-tệp.png}{caption}
% Nếu tệp chưa có, in khung nhắc để báo cáo vẫn biên dịch được.
\newcommand{\anh}[3][0.95]{%
\begin{figure}[H]
\centering
\IfFileExists{#2}{%
  \includegraphics[width=#1\textwidth]{#2}%
}{%
  \fbox{\begin{minipage}[c][4.2cm][c]{0.9\textwidth}
  \centering
  {\color{codegray}\small\textbf{[CHƯA CÓ ẢNH]}}\\[4pt]
  {\small Đặt tệp \texttt{#2} cùng thư mục với \texttt{.tex}}
  \end{minipage}}%
}
\caption{#3}
\end{figure}}

\begin{document}

% =====================================================================
%  PHẦN 1 -- TRANG BÌA & THÔNG TIN NHÓM
% =====================================================================
\begin{titlepage}
\thispagestyle{empty}
\begin{center}
{\large\bfseries TRƯỜNG CÔNG NGHỆ THÔNG TIN -- ĐẠI HỌC PHENIKAA}\\[2pt]
{\normalsize Khoa Kỹ thuật Phần mềm}\\[28pt]

{\Large\bfseries BÁO CÁO ĐỒ ÁN MÔN HỌC}\\[8pt]
{\large CSE703153 -- Kỹ thuật Lập trình An toàn (Software Security)}\\[28pt]

{\normalsize Đề tài:}\\[6pt]
{\Large\bfseries\color{primary} XÂY DỰNG VÀ KIỂM CHỨNG AN TOÀN\\[3pt]
HỆ THỐNG XỬ LÝ ĐƠN HÀNG VÀ THANH TOÁN TRỰC TUYẾN}\\[26pt]

\begin{tabular}{rl}
\textbf{Nhóm:} & 9 \\
\textbf{Lớp:} & \mustfill{Điền mã lớp} \\
\textbf{Thành viên:} & Phạm Văn Hoàng -- 23010245 \\
\textbf{Giảng viên hướng dẫn:} & MSc. Vũ Quang Dũng \\
\end{tabular}

\vfill
{\normalsize Hà Nội, tháng 10 năm 2026}
\end{center}
\end{titlepage}

\pagenumbering{arabic}
\tableofcontents
\vspace{8pt}
\noindent\emph{Ghi chú: trang bìa là Phần 1; các mục 1--7 dưới đây tương ứng
Phần 2--8 của cấu trúc bắt buộc trong hướng dẫn báo cáo.}
\clearpage

% =====================================================================
%  PHẦN 2 -- TỔNG QUAN ĐỀ TÀI
% =====================================================================
\section{Tổng quan đề tài}

\subsection{Bối cảnh và bài toán}
Hệ thống thương mại điện tử xử lý ba miền dữ liệu có quan hệ chặt chẽ: đơn hàng,
thanh toán và tồn kho. Một lỗi ở bất kỳ miền nào cũng có thể dẫn đến hậu quả
nghiệp vụ không thể hoàn tác: khách bị trừ tiền hai lần, kho bị trừ nhưng đơn
không được tạo, đơn ở trạng thái ``đã hủy'' nhưng vẫn thanh toán được, hoặc sau
khi tiến trình API chết thì hệ thống không biết giao dịch trước đó đã thành công
hay chưa. Báo cáo này trình bày một hệ thống đơn hàng -- thanh toán viết bằng
ASP.NET Core, dùng chung PostgreSQL, và quan trọng hơn là \emph{cách kiểm chứng}
các bất biến an toàn của nó bằng mô hình, SMT, phân tích bộ nhớ, fuzzing và kiểm
thử phục hồi. Nội dung tương ứng các tuần của học phần về mô hình hóa và kiểm
chứng, an toàn bộ nhớ, phân tích tĩnh/SMT và fuzzing.

\subsection{Chức năng của hệ thống}
Hệ thống hỗ trợ đăng ký/đăng nhập bằng JWT, quản lý sản phẩm và tồn kho, giỏ
hàng, tạo đơn với giá lấy từ catalog phía máy chủ, thanh toán qua một gateway mô
phỏng cục bộ có journal bền vững, hoàn tiền, và tra cứu lịch sử đơn. Các module
được tách theo miền (User, Product, Cart, Order, Payment) nhưng chạy trong cùng
một host ASP.NET Core và cùng một cơ sở dữ liệu. Mỗi thao tác nhạy cảm đều đi qua
middleware xác thực, giới hạn tần suất, giới hạn kích thước body và ghi audit.

\subsection{Kiến trúc}
Hình~\ref{fig:arch} mô tả kiến trúc tổng thể. Điểm đáng chú ý là gateway thanh
toán nằm ngoài transaction cơ sở dữ liệu: trạng thái \texttt{PROCESSING} và khóa
idempotency được ghi trước khi gọi gateway; kết quả chỉ được chốt trong một
transaction riêng. Khi kết quả không rõ, hệ thống giữ trạng thái \texttt{UNKNOWN}
và reservation tồn kho, rồi worker đối soát lại bằng cùng khóa thay vì charge mù.

\begin{figure}[H]
\centering
\begin{tikzpicture}[
  font=\small,
  box/.style={draw, rounded corners=2pt, align=center, inner sep=4pt, minimum height=9mm},
  flow/.style={-{Latex[length=2mm]}, thick},
  note/.style={font=\scriptsize\itshape, text=codegray}
]
\node[box, fill=gray!8, text width=4.2cm] (client) at (0,0)
  {Shop SPA\\\scriptsize\texttt{index.html} + \texttt{app.js}};
\node[box, fill=blue!5, text width=8.8cm] (api) at (0,-2.1)
  {\textbf{ASP.NET Core host}\\\scriptsize
   Middleware: JWT, admission/rate limit, body limit, audit\\
   Controllers $\rightarrow$ Services: User, Product, Cart, Order, Payment\\
   Validators + Order/Payment state machines};
\node[box, fill=orange!8, text width=4.0cm] (db) at (-4.9,-4.5)
  {PostgreSQL\\\scriptsize EF Core, migrations, unique/CHECK};
\node[box, fill=orange!8, text width=4.3cm] (gw) at (0,-4.5)
  {DurableFake gateway\\\scriptsize journal file + lock};
\node[box, fill=orange!8, text width=4.0cm] (worker) at (4.9,-4.5)
  {Payment recovery worker\\\scriptsize reconcile theo khóa};
\node[box, dashed, fill=green!4, text width=11.6cm] (tools) at (0,-6.5)
  {\scriptsize\textbf{Kiểm chứng:} xUnit (149 + 35) $\cdot$ Python model + TLC
   $\cdot$ Z3 SMT $\cdot$ CBMC/Valgrind/sanitizer $\cdot$ fuzz .NET/HTTP/browser
   $\cdot$ kịch bản PostgreSQL stop/start};
\draw[flow] (client) -- (api);
\draw[flow] (api.south -| db.north) -- (db);
\draw[flow] (api.south) -- (gw);
\draw[flow] (api.south -| worker.north) -- (worker);
\draw[flow, dashed] (worker.north) -- ++(0,0.8) -| (gw.east);
\draw[flow, dashed, color=green!50!black] (tools.north) -- ++(0,0.6) -| (api.south east);
\end{tikzpicture}
\caption{Kiến trúc hệ thống và các lớp kiểm chứng. Mũi tên nét đứt thể hiện luồng
đối soát/kiểm chứng, không phải luồng request đồng bộ. Khối xanh nét đứt là phần
nhóm bổ sung so với v0.}
\label{fig:arch}
\end{figure}

\subsection{Phạm vi đã thực hiện và không thực hiện}
\begin{table}[H]
\centering
\small
\renewcommand{\arraystretch}{1.25}
\begin{tabularx}{\textwidth}{@{}L{4.6cm} Y@{}}
\toprule
\textbf{Đã thực hiện} & \textbf{Chưa thực hiện / ngoài phạm vi} \\
\midrule
Mô hình threat, 27 yêu cầu, 18 trạng thái, 10 property CTL, 50 testcase đặc tả &
Không chứng minh mọi interleaving, failover, network partition hoặc power loss \\
Kiểm chứng 5 lỗi bộ nhớ C (22 ca động, CBMC 10 ca, Valgrind 10 ca, SAT/DRAT 10 ca) &
Các lỗi C chỉ là mô phỏng didactic, không phải lỗi trong runtime C\# \\
\addlinespace
6 CWE thiết kế/nghiệp vụ đã vá: 190, 841, 306, 602, 20, 639 (unit/integration/HTTP) &
Chưa kiểm mọi route mới sinh về sau; gateway thật ngoài phạm vi \\
SMT cho số học thanh toán (7 ca) và tồn kho (3 ca) &
Không mô hình hóa toàn bộ API, EF Core hay gateway thật \\
TLC cho \texttt{OrderPayment.tla} (35 trạng thái, depth 9) &
Chưa có chứng minh vô hạn cho C\#; TLC chỉ đúng với cấu hình hữu hạn \\
149 unit + 35 integration test, fuzz có seed cố định, fuzz HTTP/browser &
Chưa code-coverage-guided/path/taint cho app ASP.NET \\
Phục hồi charge/refund sau khi PostgreSQL stop/start và sau khi kill API &
Chưa thử provider thật, TLS/proxy, đa vùng hoặc Kubernetes \\
\bottomrule
\end{tabularx}
\caption{Phạm vi công việc. Các giới hạn này được nhắc lại trong Phần 6.}
\end{table}

\subsection{Công nghệ sử dụng}
\begin{table}[H]
\centering
\small
\renewcommand{\arraystretch}{1.2}
\begin{tabularx}{\textwidth}{@{}L{4.2cm} Y@{}}
\toprule
\textbf{Thành phần} & \textbf{Công nghệ / phiên bản} \\
\midrule
Ứng dụng & .NET 10, ASP.NET Core, EF Core 8, Npgsql \\
Cơ sở dữ liệu & PostgreSQL 16.x (Testcontainers.PostgreSql 4.15.0), migration + unique/CHECK constraints \\
Kiểm thử & xUnit, Coverlet, WebApplicationFactory, Testcontainers/PostgreSQL local \\
Mô hình & Python 3, JSON state model, TLA+ 1.7.4 (TLC) \\
SMT/BMC & Z3, CBMC (unwind 64), CaDiCaL + drat-trim \\
Bộ nhớ & GCC ASan/UBSan/LSan, Valgrind Memcheck, Cppcheck \\
Fuzzing & Fuzz seed cố định, HTTP black-box, Chrome CDP, trace-pc C fixture \\
\bottomrule
\end{tabularx}
\caption{Công nghệ và công cụ chính.}
\end{table}

% =====================================================================
%  PHẦN 3 -- PHƯƠNG PHÁP LUẬN
% =====================================================================
\section{Phương pháp luận}

\subsection{Quy trình năm bước}
Nhóm áp dụng một quy trình lặp, mỗi bước sinh ra một artifact kiểm chứng được:

\begin{enumerate}
\item \textbf{Đặc tả và threat model.} Từ luồng nghiệp vụ, nhóm ghi lại tài sản,
ranh giới tin cậy, threat theo STRIDE, yêu cầu bảo mật và các bất biến
(đơn đã hủy không thể thanh toán, tổng tiền khớp, kho không âm, refund chỉ một lần).
Kết quả là workbook 10 sheet và các model JSON.
\item \textbf{Model checking.} Chuyển bất biến thành property trên Kripke structure
và TLA+; kiểm bằng checker độc lập và TLC. Nhóm cố ý tạo các mutation âm để chứng
minh checker thực sự bắt được lỗi, thay vì chỉ in ``PASS''.
\item \textbf{An toàn bộ nhớ.} Mô phỏng 5 lớp lỗi C (rò rỉ, double free,
use-after-free, tràn heap, tràn stack) trong các module cùng tên với miền nghiệp vụ;
so sánh bản BUGGY và FIXED bằng ASan/UBSan/LSan, Cppcheck, CBMC và Valgrind.
\item \textbf{SMT và BMC.} Mã hóa số học thanh toán/tồn kho thành công thức Z3 để
tìm counterexample hoặc chứng minh UNSAT trong miền đã chọn; dùng CBMC cho các
harness C với bound xác định.
\item \textbf{Kiểm chứng động và phục hồi.} Chạy unit/integration test thật,
fuzz theo seed cố định, fuzz HTTP/browser, rồi mô phỏng crash: dừng PostgreSQL và
kill tiến trình API giữa charge/refund để kiểm tra tính idempotent của recovery.
\end{enumerate}

\subsection{Công cụ và cấu hình}
\begin{table}[H]
\centering
\footnotesize
\renewcommand{\arraystretch}{1.2}
\begin{tabularx}{\textwidth}{@{}L{3.4cm} L{5.4cm} Y@{}}
\toprule
\textbf{Công cụ} & \textbf{Cấu hình / cờ} & \textbf{Mục đích} \\
\midrule
GCC & \cmd{-std=c11 -Wall -Wextra -Wpedantic}, ASan/UBSan/LSan & Biên dịch và phát hiện lỗi bộ nhớ \\
Cppcheck & \cmd{--enable=all --inconclusive} & Phân tích tĩnh độc lập \\
CBMC & \cmd{--bounds-check --pointer-check --memory-leak-check --unwind 64} & BMC có bound cho 5 module C \\
Valgrind & Memcheck, build không sanitizer & Kiểm tra bộ nhớ độc lập với ASan \\
Z3 & QF\_NIA, integer cents & SMT cho số học thanh toán/tồn kho \\
TLC & \cmd{MaxAttempts=2}, 1 worker & Model checking TLA+ hữu hạn \\
Coverlet & cobertura, include 5 target & Coverage dòng/nhánh \\
Chrome CDP & headless, chặn response & Fuzz lỗi mạng phía UI \\
\bottomrule
\end{tabularx}
\caption{Công cụ và cấu hình đã dùng.}
\end{table}

\subsection{Sinh dữ liệu và chính sách seed}
Mọi campaign đều dùng seed cố định để tái lập: 20261009 cho fuzz .NET,
20261022 cho fuzz HTTP và C coverage-guided. Fuzz .NET chạy 25.000
iteration cho campaign validator (order + payment) và 2.000 chuỗi trạng thái,
mỗi chuỗi 20 bước. Fuzz HTTP
gửi 500 mutation (tổng 601 request) lên API thật với PostgreSQL tạm; fuzz C chạy
20.000 mutation trên fixture tổng hợp. Không có bước \cmd{pip install} bắt buộc;
các script chỉ dùng thư viện chuẩn của Python. Khác với các đề tài có sẵn script
\texttt{data\_gen/}, dự án này sinh dữ liệu fuzz từ corpus seed và tham số
\texttt{-{}-seed}/\texttt{-{}-iterations} như bảng campaign ở Phần 5.

\subsection{Harness tự viết}
Ngoài các test có sẵn, nhóm viết thêm bốn harness: (i) checker Kripke/liveness bằng
Python duyệt toàn bộ state space và sinh counterexample; (ii) harness CBMC cho 5
module C, include trực tiếp module thật; (iii) HTTP fuzzer khởi tạo PostgreSQL tạm,
chạy API thật và kiểm bất biến DB sau mỗi request; (iv) kịch bản recovery dừng
PostgreSQL/API rồi khởi động lại, đối chiếu journal và số lần gọi gateway.

\subsection{Ánh xạ v0/v1 trong repository}
Repo dùng hai biến thể BUGGY và FIXED trên cùng một file C, chọn
bằng \cmd{-UFIXED} và \cmd{-DFIXED}, thay vì hai file \texttt{*\_v0.c}/\texttt{*\_v1.c}.
Trong báo cáo, \texttt{v0} tương ứng BUGGY và \texttt{v1} tương ứng FIXED; mọi
trích đoạn đều ghi rõ dòng gốc. Với các lỗi tầng thiết kế/nghiệp vụ
(CWE-190, CWE-841, CWE-306, CWE-602, CWE-20, CWE-639), v0 là model ``buggy'' hoặc
tình trạng trước remediation ghi trong \path{security-analysis/vulnerabilities.json},
còn v1 là mã hiện tại đã có guard/validator/authorization tương ứng. Riêng
CWE-306/602/20/639, đoạn v0 trong báo cáo là minh họa ngắn tái dựng theo finding
(không có file v0 trong repo).

% =====================================================================
%  PHẦN 4 -- PHÂN TÍCH LỖ HỔNG & BẢN VÁ
% =====================================================================
\section{Phân tích lỗ hổng \& Bản vá}

Phần này trình bày 11 lỗ hổng theo đúng 5 mục con: vị trí, nguyên nhân gốc rễ,
bằng chứng lỗi tồn tại, cách vá (trước/sau) và bằng chứng đã vá. Bảy mục đầu gồm
năm module C mô phỏng trực tiếp thao tác nguy hiểm của miền nghiệp vụ và hai lỗi
thiết kế ở tầng đơn hàng/thanh toán (CWE-190, CWE-841) kiểm chứng bằng model/SMT.
Bốn mục cuối (CWE-306, CWE-602, CWE-20, CWE-639) là finding bảo mật tầng
API/nghiệp vụ, kiểm chứng bằng unit, integration và HTTP. Các listing C lược bớt
chỉ thị \texttt{\#ifdef}/\texttt{\#endif} để dễ đối chiếu trước/sau; số dòng vẫn giữ
theo file gốc và ghi rõ biến thể.

% ---------------------------------------------------------------------
\subsection{\cwe{401} -- Rò rỉ bộ nhớ khi hủy đơn}
\paragraph{Vị trí và mô tả.}
\path{security-analysis/c-memory/modules/order.c}, hàm \cmd{cancel_order},
dòng 6--18. Hàm cấp phát \cmd{Order} và mảng \cmd{OrderItem} lồng bên trong, sau đó
giải phóng đối tượng cha.

\paragraph{Nguyên nhân gốc rễ.}
C không có cơ chế thu gom rác. Lời gọi \cmd{free(order)} chỉ trả vùng nhớ của
struct cha, không giải phóng vùng \cmd{order->items} được cấp phát ở dòng 10.
Mỗi lần hủy đơn trong vòng lặp 5 lần lại rò rỉ một mảng, nên LeakSanitizer phát
hiện leak có kiểm soát. Đây là \cwe{401} (Missing Release of Memory after Effective
Lifetime), tương ứng quy tắc CERT MEM51-C.

\paragraph{Bằng chứng lỗi tồn tại.}
Biên dịch v0 bằng \cmd{-UFIXED} và chạy, chương trình thoát với mã 1; evidence ghi
nhận \cmd{LeakSanitizer: detected memory leaks}. Log đầy đủ nằm ở
\path{security-analysis/c-memory/results/order_buggy.log} (sinh bởi
\cmd{make -C security-analysis/c-memory test}).

\begin{lstlisting}[language=C,numbers=none,caption={v0 -- trích \texttt{order.c}, dòng 6--18 (BUGGY, \texttt{-UFIXED})},firstnumber=6]
static void cancel_order(void) {
    Order *order = allocate(sizeof(*order));
    order->id = 1;
    order->count = 3;
    order->items = allocate(order->count * sizeof(*order->items));
    for (size_t i = 0; i < order->count; ++i)
        order->items[i] = (OrderItem){(int)i + 1, 2, 45000};
    printf("Cancel order %d: %zu items\n", order->id, order->count);
    free(order);          /* chi giai phong struct cha */
}
\end{lstlisting}

\paragraph{Cách vá.}
Giải phóng mảng con trước khi giải phóng cha. Bản v1 thêm dòng \cmd{free(order->items)}
ngay trước \cmd{free(order)}; đây là sửa nguyên nhân gốc rễ chứ không chỉ làm im
triệu chứng.

\begin{lstlisting}[language=C,numbers=none,caption={v1 -- trích \texttt{order.c}, dòng 14--17 (FIXED)},firstnumber=14]
    free(order->items);   /* v1: giai phong mang con truoc */
    free(order);          /* ca v0 va v1 */
\end{lstlisting}

\paragraph{Bằng chứng đã vá.}
Bản v1 chạy sạch: exit code 0, không có báo cáo sanitizer; CBMC kết luận
\texttt{VERIFICATION SUCCESSFUL} trong bound 64 và Valgrind không phát hiện leak.
Bảng~\ref{tab:beforeafter} tóm tắt trạng thái trước/sau.

% ---------------------------------------------------------------------
\subsection{\cwe{415} -- Double free khi dọn dẹp giao dịch lỗi}
\paragraph{Vị trí và mô tả.}
\path{security-analysis/c-memory/modules/payment.c}, hàm \cmd{main}, dòng 5--14.
Nhánh lỗi gateway giải phóng \cmd{transaction_code}, sau đó phần cleanup chung
giải phóng lại cùng con trỏ.

\paragraph{Nguyên nhân gốc rễ.}
Hai đoạn mã cùng tin rằng mình sở hữu vùng nhớ. Sau lần \cmd{free} đầu tiên
(dòng 9), con trỏ không được đặt về \cmd{NULL}, nên lần \cmd{free} thứ hai
(dòng 13) giải phóng cùng một địa chỉ. Đây là \cwe{415} (Double Free), quy tắc
CERT MEM31-C.

\paragraph{Bằng chứng lỗi tồn tại.}
ASan báo \cmd{AddressSanitizer: attempting double-free} và tiến trình dừng với mã
khác 0. Log đầy đủ ở \path{security-analysis/c-memory/results/payment_buggy.log}.

\begin{lstlisting}[language=C,numbers=none,caption={v0 -- trích \texttt{payment.c}, dòng 6--13 (BUGGY)},firstnumber=6]
Payment payment = {1, allocate(32)};
snprintf(payment.transaction_code, 32, "%s", "TX-DEMO-001");
puts("Gateway failure: release transaction code");
free(payment.transaction_code);      /* lan 1 */
free(payment.transaction_code);      /* BUG: lan 2, cung dia chi */
\end{lstlisting}

\paragraph{Cách vá.}
Sau khi giải phóng, gán con trỏ về \cmd{NULL}. Lần \cmd{free(NULL)} tiếp theo là
hợp lệ và không giải phóng lần hai; cách này loại bỏ trạng thái ``con trỏ treo''
chứ không chỉ che lỗi.

\begin{lstlisting}[language=C,numbers=none,caption={v1 -- trích \texttt{payment.c}, dòng 9--13 (FIXED)},firstnumber=9]
free(payment.transaction_code);      /* lan 1 */
payment.transaction_code = NULL;     /* v1: khong con tro treo */
free(payment.transaction_code);      /* free(NULL) an toan */
\end{lstlisting}

\paragraph{Bằng chứng đã vá.}
Bản FIXED exit 0, ASan/Valgrind sạch; CBMC kết luận
\texttt{VERIFICATION SUCCESSFUL} cho harness trong bound 64. Trong bảng tổng hợp,
đây là ca duy nhất có lỗi ``attempting double-free''.

% ---------------------------------------------------------------------
\subsection{\cwe{416} -- Use-after-free khi ghi audit xóa giỏ hàng}
\paragraph{Vị trí và mô tả.}
\path{security-analysis/c-memory/modules/cart.c}, hàm \cmd{main}, dòng 5--19.
Sau khi xóa một \cmd{CartItem}, đoạn log lại đọc \cmd{removed->product_id} từ
vùng nhớ đã giải phóng.

\paragraph{Nguyên nhân gốc rễ.}
Thứ tự thao tác sai: \cmd{free(removed)} xảy ra trước khi đọc dữ liệu cần log.
Sau \cmd{free}, vùng nhớ không còn thuộc về chương trình; đọc lại là hành vi
không xác định. Đây là \cwe{416} (Use After Free), CERT MEM30-C.

\paragraph{Bằng chứng lỗi tồn tại.}
ASan báo \cmd{AddressSanitizer: heap-use-after-free}; exit code khác 0. Log đầy đủ
ở \path{security-analysis/c-memory/results/cart_buggy.log}.

\begin{lstlisting}[language=C,numbers=none,caption={v0 -- trích \texttt{cart.c}, dòng 16--17 (BUGGY)},firstnumber=16]
free(removed);
printf("Removed product %d\n", removed->product_id); /* BUG: doc sau free */
\end{lstlisting}

\paragraph{Cách vá.}
Sao chép dữ liệu cần thiết (\cmd{product_id}) ra biến cục bộ \emph{trước} khi
giải phóng, rồi log từ bản sao; sau đó đặt con trỏ về \cmd{NULL}. Cách này loại bỏ
việc truy cập vùng nhớ đã free.

\begin{lstlisting}[language=C,numbers=none,caption={v1 -- trích \texttt{cart.c}, dòng 10--14 (FIXED)},firstnumber=10]
int product_id = removed->product_id; /* copy truoc khi free */
free(removed);
removed = NULL;
printf("Removed product %d\n", product_id);
\end{lstlisting}

\paragraph{Bằng chứng đã vá.}
Bản FIXED exit 0, ASan/Valgrind sạch; CBMC kết luận
\texttt{VERIFICATION SUCCESSFUL} cho harness trong bound 64.

% ---------------------------------------------------------------------
\subsection{\cwe{122} -- Tràn heap khi nhập tên sản phẩm}
\paragraph{Vị trí và mô tả.}
\path{security-analysis/c-memory/modules/product.c}, hàm \cmd{main}, dòng 6--20.
Tên sản phẩm được sao chép vào buffer heap 32 byte bằng \cmd{strcpy}.

\paragraph{Nguyên nhân gốc rễ.}
\cmd{strcpy} không nhận kích thước đích và không kiểm tra độ dài nguồn. Với input
dài hơn 32 byte, nó ghi qua biên allocation. Đây là \cwe{122} (Heap-based Buffer
Overflow); quy tắc CERT STR31-C yêu cầu bảo đảm đủ chỗ cho dữ liệu và ký tự kết thúc.

\paragraph{Bằng chứng lỗi tồn tại.}
ASan báo \cmd{AddressSanitizer: heap-buffer-overflow} khi chạy với input mặc định
dài; log ở \path{security-analysis/c-memory/results/product_buggy.log}. Cppcheck
không bắt được lỗi này (chỉ bắt lỗi biên của \texttt{user.c}), cho thấy cần nhiều
công cụ bổ trợ.

\begin{lstlisting}[language=C,numbers=none,caption={v0 -- trích \texttt{product.c}, dòng 6--17},firstnumber=6]
Product product = {101, allocate(NAME_CAPACITY)};
...
strcpy(product.name, input);   /* khong kiem tra do dai */
\end{lstlisting}

\paragraph{Cách vá.}
Kiểm tra \cmd{strlen(input) >= NAME_CAPACITY} trước khi copy; nếu vượt, từ chối
input và giải phóng buffer. Bản vá chọn \emph{từ chối} thay vì cắt bớt âm thầm,
để không tạo dữ liệu sai lệch.

\begin{lstlisting}[language=C,numbers=none,caption={v1 -- trích \texttt{product.c}, dòng 10--17 (FIXED)},firstnumber=10]
if (strlen(input) >= NAME_CAPACITY) {
    puts("REJECTED: product name too long");
    free(product.name);
    return 0;
}
strcpy(product.name, input);  /* an toan sau khi kiem tra */
\end{lstlisting}

\paragraph{Bằng chứng đã vá.}
Bản FIXED sạch với các độ dài biên 0, 1, 31, 32 và 128 (5 ca bổ sung trong bộ 22
ca); ASan/CBMC/Valgrind đều kết luận an toàn trong phạm vi đã kiểm.

% ---------------------------------------------------------------------
\subsection{\cwe{121} -- Tràn stack khi đăng ký email}
\paragraph{Vị trí và mô tả.}
\path{security-analysis/c-memory/modules/user.c}, hàm \cmd{register_user},
dòng 5--15. Email được copy vào mảng cục bộ \cmd{char email[32]}.

\paragraph{Nguyên nhân gốc rễ.}
Cùng bản chất \cwe{122} nhưng vùng đích nằm trên stack. Ghi quá 32 byte sẽ ghi đè
biến cục bộ, return address hoặc con trỏ frame, gây crash hoặc hành vi không xác
định. Đây là \cwe{121} (Stack-based Buffer Overflow), CERT STR31-C.

\paragraph{Bằng chứng lỗi tồn tại.}
ASan báo \cmd{AddressSanitizer: stack-buffer-overflow}; với input đủ dài,
Valgrind có thể báo invalid jump do stack corruption, còn CBMC chỉ ra chính xác
ghi ngoài \cmd{email[32]}. Log ở \path{security-analysis/c-memory/results/user_buggy.log}.

\begin{lstlisting}[language=C,numbers=none,caption={v0 -- trích \texttt{user.c}, dòng 5--15},firstnumber=5]
static void register_user(const char *input) {
    char email[EMAIL_CAPACITY];
    strcpy(email, input);   /* khong kiem tra do dai */
    printf("Registered email: %s\n", email);
}
\end{lstlisting}

\paragraph{Cách vá.}
Kiểm tra độ dài trước khi copy và từ chối input quá dài. Bản v1 dùng
\cmd{strlen(input) >= sizeof(email)}; cách khác tương đương là \cmd{snprintf} với
kích thước đích.

\begin{lstlisting}[language=C,numbers=none,caption={v1 -- trích \texttt{user.c}, dòng 7--13 (FIXED)},firstnumber=7]
if (strlen(input) >= sizeof(email)) {
    puts("REJECTED: email too long");
    return;
}
strcpy(email, input);   /* an toan sau khi kiem tra */
\end{lstlisting}

\paragraph{Bằng chứng đã vá.}
Bản FIXED sạch ở 6 ca độ dài (0, 1, 31, 32, 128 và input hợp lệ); CBMC
\texttt{VERIFICATION SUCCESSFUL}, Valgrind không còn lỗi.

% ---------------------------------------------------------------------
\subsection{\cwe{190} -- Tràn số nguyên trong tính tiền thanh toán}
\paragraph{Vị trí và mô tả.}
Mô hình \path{OrderPaymentMicroservice/Security/PaymentVerification.z3},
các test \cmd{TEST6_INT32_SUBTOTAL_OVERFLOW_EXISTS} (dòng 93--100) và
\cmd{TEST7_TOTAL_BOUND_BLOCKS_OVERSIZED_SUBTOTAL} (dòng 102--109). Phía
ứng dụng, ràng buộc tương ứng nằm ở \path{OrderPaymentMicroservice/Security/PaymentValidator.cs}
dòng 21--46 và \path{OrderPaymentMicroservice/Security/MoneyAmountAttribute.cs} dòng 9--13.

\paragraph{Nguyên nhân gốc rễ.}
Nếu subtotal được tính bằng số nguyên 32-bit có dấu, tích
\cmd{quantity * unitPriceCents} có thể vượt $2^{31}-1$ ngay cả khi từng trường
nằm trong biên hợp lệ: $10\,000 \times 99\,999\,999 = 999\,999\,990\,000$.
Kết quả tràn thành số âm hoặc sai, làm lệch tổng đơn và số tiền thanh toán.
Đây là \cwe{190} (Integer Overflow or Wraparound), quy tắc CERT INT32-C.

\paragraph{Bằng chứng lỗi tồn tại.}
Z3 trả \texttt{sat} cho TEST6: tồn tại cặp \cmd{Quantity1},
\cmd{UnitPrice1Cents} hợp lệ theo \cmd{ValidItem} nhưng tích vượt int32. Đây là
counterexample cho giả định ``biên từng trường là đủ''.

\begin{lstlisting}[language=C,numbers=none,caption={v0 -- mô hình int32 trong \texttt{PaymentVerification.z3}, dòng 93--100},firstnumber=93]
; Counterexample: allowed item bounds can overflow a signed 32-bit subtotal.
(echo "TEST6_INT32_SUBTOTAL_OVERFLOW_EXISTS")
(push)
(assert (ValidItem Quantity1 UnitPrice1Cents))
(assert (> (* Quantity1 UnitPrice1Cents) 2147483647))
(check-sat)
(get-value (Quantity1 UnitPrice1Cents (* Quantity1 UnitPrice1Cents)))
(pop)
\end{lstlisting}

\paragraph{Cách vá.}
Ba lớp bảo vệ: (i) dùng \cmd{decimal} thay vì \cmd{int} cho tiền; (ii) chặn tổng
đơn ở \cmd{999999.99} (tương đương 99.999.999 cent); (iii) bắt buộc tiền là bội số
của 0,01 thay vì làm tròn âm thầm. Với các ràng buộc này, TEST7 là \texttt{unsat}:
không thể có subtotal vượt biên trong một đơn hợp lệ.

\begin{lstlisting}[language=CSharp,numbers=none,caption={v1 -- ràng buộc trong \texttt{PaymentValidator.cs}, dòng 37--40},firstnumber=37]
if (request.Amount > MAX_PAYMENT_AMOUNT)
    throw new ArgumentException($"Amount cannot exceed {MAX_PAYMENT_AMOUNT}", nameof(request.Amount));

MoneyAmounts.RequireCents(request.Amount, nameof(request.Amount));
\end{lstlisting}

\paragraph{Bằng chứng đã vá.}
Bộ 7 ca SMT đều đúng kỳ vọng (7/7), trong đó TEST7 \texttt{unsat}; validator và
helper đạt branch coverage 97,22\% và 100\%. Giới hạn: mô hình dùng số nguyên toán
học, chưa chứng minh mọi đường thực thi C\#.

% ---------------------------------------------------------------------
\subsection{\cwe{841} -- Đơn đã hủy vẫn có thể bị thanh toán}
\paragraph{Vị trí và mô tả.}
Model \path{security-analysis/order_state_model_buggy.json} (9 trạng thái,
16 chuyển tiếp) mô tả thiết kế v0; witness
\cmd{CREATED -> CANCELLED -> PAID_AFTER_CANCELLED}. Phía mã hiện tại, guard nằm
ở \path{OrderPaymentMicroservice/Services/PaymentService.cs} dòng 55--56 và bảng
chuyển tiếp ở \path{OrderPaymentMicroservice/Security/OrderStateMachine.cs} dòng 10--22.

\paragraph{Nguyên nhân gốc rễ.}
Thiếu kiểm tra trạng thái đơn trước khi tạo thanh toán. Nếu chỉ dựa vào
\cmd{OrderId}, một request đến sau khi đơn đã hủy vẫn tạo được payment; đây là
\cwe{841} (Improper Enforcement of Behavioral Workflow), không phải lỗi bộ nhớ.
Lỗi tương tự cũng có thể xuất hiện dưới dạng double charge nếu hai request đồng
thời cùng vượt qua bước kiểm tra. Hình~\ref{fig:order-state} trực quan hóa bảng
chuyển tiếp và chuyển tiếp bất hợp lệ.

\begin{figure}[H]
\centering
\IfFileExists{kripke-order.pdf}{%
  \includegraphics[width=0.95\textwidth]{kripke-order.pdf}%
}{%
  \IfFileExists{kripke-order.png}{%
    \includegraphics[width=0.95\textwidth]{kripke-order.png}%
  }{%
    \fbox{\begin{minipage}[c][4.2cm][c]{0.9\textwidth}\centering
    {\color{codegray}\small\textbf{[CHƯA CÓ ẢNH KRIPKE]}}\\[4pt]
    {\small Đặt tệp \texttt{kripke-order.pdf} (hoặc \texttt{.png}) cùng thư mục với \texttt{.tex}}
    \end{minipage}}%
  }%
}
\caption{Trừu tượng hóa máy trạng thái Order (rút từ \texttt{OrderStateMachine.cs})
dưới dạng cấu trúc Kripke: nút là trạng thái, nhãn nút là tập mệnh đề nguyên tử,
cạnh là chuyển tiếp. Mô hình v0 cho phép chuyển tiếp bất hợp lệ
\texttt{CANCELLED -> PAID}; bản v1 chặn bằng guard trạng thái nguồn.}
\label{fig:order-state}
\end{figure}

\paragraph{Bằng chứng lỗi tồn tại.}
Checker trả \cmd{P1 holds=false} với witness states/events ghi rõ
\cmd{"CreatePaymentAsync: missing order status guard"}. Kết quả nằm ở
\path{security-analysis/results/order_state_model_buggy.results.json}.

\begin{lstlisting}[language=json,caption={v0 -- witness từ model buggy},firstnumber=1]
"id": "P1",
"formula": "invalid_transition",
"holds": false,
"witness_states": ["CREATED", "CANCELLED", "PAID_AFTER_CANCELLED"],
"witness_events": ["cancel", "CreatePaymentAsync: missing order status guard"]
\end{lstlisting}

\paragraph{Cách vá.}
Đưa kiểm tra trạng thái vào đúng nơi tạo payment, đồng thời chuẩn hóa bảng chuyển
tiếp để mọi đường đi đều phải qua \cmd{OrderStateMachine}. Thêm khóa idempotency,
unique constraint và transaction \cmd{Serializable} để chặn hai request đồng thời.

\begin{lstlisting}[language=CSharp,numbers=left,caption={v1 -- guard trong \texttt{PaymentService.cs}, dòng 44--60},firstnumber=44]
var previousAttempt = order.Payments.FirstOrDefault(candidate =>
    candidate.IdempotencyKey == request.IdempotencyKey);
if (previousAttempt != null)
{
    if (previousAttempt.Amount != request.Amount)
        throw new InvalidOperationException("Idempotency key was already used with a different amount");
    return MapToResponse(previousAttempt);
}
if (order.Payments.Count != 0)
    throw new InvalidOperationException("Order already owns a logical payment; retry that payment instead");
if (order.Status != OrderStatus.PAYMENT_PENDING)
    throw new InvalidOperationException($"Order in state {order.Status} cannot be paid");
if (order.InventoryStatus != InventoryReservationStatus.RESERVED)
    throw new InvalidOperationException("Order has no active inventory reservation");
if (request.Amount != order.TotalAmount)
    throw new InvalidOperationException("Payment amount must equal the order total");
\end{lstlisting}

\paragraph{Bằng chứng đã vá.}
Model đúng đạt 15/15 property, 6/6 mutation âm bị phát hiện; TLC kiểm 5 invariant +
termination trên 35 trạng thái; 35 integration test (gồm idempotency, audit nguyên
tử, kill API) đều PASS. Phần 5 trình bày thêm kịch bản recovery.

% ---------------------------------------------------------------------
\subsection{\cwe{306} -- Truy cập Order khi chưa xác thực}
\paragraph{Vị trí và mô tả.}
\path{OrderPaymentMicroservice/Controllers/OrdersController.cs} (thuộc tính cấp lớp)
và \path{OrderPaymentMicroservice/Program.cs} dòng 73--88 (cấu hình JWT). Finding
\cmd{VULN-001}.

\paragraph{Nguyên nhân gốc rễ.}
Giai đoạn v0 thiếu middleware xác thực và \cmd{[Authorize]}, cũng không kiểm quyền
sở hữu theo claim. Client đọc, liệt kê và hủy đơn của bất kỳ ai mà không cần chứng
minh danh tính; đây là \cwe{306} (Missing Authentication for Critical Function),
kèm khía cạnh \cwe{862}.

\paragraph{Bằng chứng lỗi tồn tại.}
Finding \cmd{VULN-001} với \cmd{detected_by = Static code review}: endpoint Order
không gắn policy quyền. Đoạn v0 dưới đây \emph{minh họa} tình trạng trước
remediation (không phải file trong repo):

\begin{lstlisting}[language=CSharp,numbers=none,breaklines=true,caption={v0 -- minh họa OrdersController thiếu xác thực (theo VULN-001)}]
[ApiController]
[Route("api/[controller]")]
public class OrdersController : ControllerBase   // v0: khong co [Authorize]
{
    [HttpGet("{id}")]
    public async Task<IActionResult> GetOrder(int id)   // v0: khong kiem ownership
        => Ok(await _orderService.GetOrderByIdAsync(id));
}
\end{lstlisting}

\paragraph{Cách vá.}
Thêm JWT authentication + \cmd{[Authorize]} ở cấp lớp, policy \cmd{Roles = "Admin"}
cho liệt kê toàn bộ, và kiểm ownership theo claim \cmd{sub} qua
\cmd{User.CanAccessUser}; \cmd{GET /api/orders/mine} luôn lấy subject thay vì nhận
ID từ client.

\begin{lstlisting}[language=CSharp,numbers=none,breaklines=true,caption={v1 -- \texttt{OrdersController.cs}, xác thực và ownership}]
[ApiController]
[Route("api/[controller]")]
[Authorize]                       // v1: bat buoc xac thuc
public class OrdersController : ControllerBase
...
if (!User.CanAccessUser(order.CustomerId))   // v1: kiem so huu
    return Forbid();
\end{lstlisting}

\paragraph{Bằng chứng đã vá.}
Integration test
\cmd{OrderPaymentApiTests.Authentication_and_resource_ownership_are_enforced}
kiểm anonymous $\to$ 401 và cross-user $\to$ 403 cho cart, order, payment,
operation; HTTP fuzz ghi 27 response 401 và 46 response 403. Giới hạn: chưa có
security test tự động bắt buộc cho mọi route sinh về sau.

% ---------------------------------------------------------------------
\subsection{\cwe{602} -- Client kiểm soát giá và tổng tiền}
\paragraph{Vị trí và mô tả.}
\path{OrderPaymentMicroservice/Services/OrderService.cs}, hàm \cmd{CreateOrderAsync}
(dòng 30--124). Finding \cmd{VULN-002}.

\paragraph{Nguyên nhân gốc rễ.}
v0 lưu thẳng \cmd{UnitPrice} và \cmd{TotalAmount} từ payload, tin dữ liệu có thẩm
quyền về giá do client cung cấp. Client có thể khai giá 1 đồng cho sản phẩm giá 25
đồng; đây là \cwe{602} (Client-Side Enforcement of Server-Side Security).

\paragraph{Bằng chứng lỗi tồn tại.}
Characterization bằng \cmd{OrderServiceTests.ClientCannotCreateOrderUsingTamperedPrice}:
gửi \cmd{TotalAmount = 2}, \cmd{UnitPrice = 1} cho sản phẩm giá 25 bị từ chối
(\cmd{ArgumentException}), không tạo order, tồn kho không đổi. Đoạn v0 minh họa:

\begin{lstlisting}[language=CSharp,numbers=none,breaklines=true,caption={v0 -- minh họa OrderService tin giá client (theo VULN-002)}]
var order = new Order {
    TotalAmount = request.TotalAmount,          // v0: tin client
    OrderItems = request.Items.Select(i => new OrderItem {
        UnitPrice = i.UnitPrice                 // v0: tin client
    }).ToList()
};
\end{lstlisting}

\paragraph{Cách vá.}
Đọc \cmd{Product.Price} từ server, tự tính \cmd{serverTotal} bằng \cmd{checked},
đối chiếu \cmd{request.TotalAmount} trong sai số 1 cent, và bọc trong transaction
\cmd{Serializable} để chặn race tồn kho.

\begin{lstlisting}[language=CSharp,numbers=none,breaklines=true,caption={v1 -- \texttt{OrderService.cs}, giá lấy từ server (dòng 73--87)}]
decimal serverTotal = 0;
foreach (var item in request.Items) {
    var product = products[item.ProductId];
    MoneyAmounts.RequireCents(product.Price, nameof(product.Price));
    if (!product.IsActive)
        throw new OrderCreationConflictException("PRODUCT_UNAVAILABLE", $"Product {product.Id} is not active");
    if (product.Stock < requestedQuantities[item.ProductId])
        throw new OrderCreationConflictException("OUT_OF_STOCK", $"Not enough stock for product {product.Id}");
    serverTotal = checked(serverTotal + item.Quantity * product.Price);
}
if (Math.Abs(request.TotalAmount - serverTotal) > 0.01m)
    throw new ArgumentException("TotalAmount does not match current product prices", nameof(request.TotalAmount));
\end{lstlisting}

\paragraph{Bằng chứng đã vá.}
\cmd{OrderServiceTests.ClientCannotCreateOrderUsingTamperedPrice} PASS; integration
flow đặt hàng hợp lệ PASS; SMT tồn kho và HTTP regression đều sạch. Cần lưu ý đây
là kiểm chứng theo mẫu hành vi (tampered price), không phải chứng minh hình thức
cho mọi định dạng payload.

% ---------------------------------------------------------------------
\subsection{\cwe{20} -- Bỏ qua lệch tổng tiền}
\paragraph{Vị trí và mô tả.}
\path{OrderPaymentMicroservice/Security/InputValidator.cs}, nhánh kiểm dung sai
trong \cmd{ValidateCreateOrderRequest} (dòng 66--72). Finding \cmd{VULN-004}.

\paragraph{Nguyên nhân gốc rễ.}
v0: điều kiện lệch quá 0,01 rơi vào một khối rỗng, nên tổng sai vẫn được chấp nhận;
đồng thời thiếu kiểm giá dưới ngưỡng tối thiểu. Đây là \cwe{20} (Improper Input
Validation).

\paragraph{Bằng chứng lỗi tồn tại.}
Unit test mô tả (characterization) chứng minh tổng 0,01 so với item 999999.99 từng
được nhận. Đoạn v0 minh họa:

\begin{lstlisting}[language=CSharp,numbers=none,breaklines=true,caption={v0 -- minh họa tolerance rỗng (theo VULN-004)}]
if (Math.Abs(request.TotalAmount - calculatedTotal) > 0.01m)
{
    // v0: khong lam gi -> mismatch bi bo qua
}
\end{lstlisting}

\paragraph{Cách vá.}
Từ chối khi lệch quá 0,01 (\cmd{throw ArgumentException}) và từ chối \cmd{UnitPrice}
dưới \cmd{MIN_PRICE} (0.01) hoặc không phải bội số cent.

\begin{lstlisting}[language=CSharp,numbers=none,breaklines=true,caption={v1 -- \texttt{InputValidator.cs} (dòng 66--101)}]
const decimal tolerance = 0.01m;
if (Math.Abs(request.TotalAmount - calculatedTotal) > tolerance)
    throw new ArgumentException("TotalAmount must match the sum of order items", nameof(request.TotalAmount));
...
if (item.UnitPrice < MIN_PRICE)
    throw new ArgumentException($"UnitPrice must be at least {MIN_PRICE}", nameof(item.UnitPrice));
MoneyAmounts.RequireCents(item.UnitPrice, nameof(item.UnitPrice));
\end{lstlisting}

\paragraph{Bằng chứng đã vá.}
\cmd{InputValidatorTests.MismatchOrBelowMinimumPriceIsRejected} (5 ca) PASS; boundary
+ fuzz oracle; branch coverage của \cmd{InputValidator} đạt 97,22\%.

% ---------------------------------------------------------------------
\subsection{\cwe{639} -- Đọc/sửa giỏ hàng của người khác}
\paragraph{Vị trí và mô tả.}
\path{OrderPaymentMicroservice/Controllers/CartController.cs}, các route
\texttt{user/\{userId\}/add} và \texttt{\{cartId\}/item/\{itemId\}}. Finding
\cmd{VULN-005}.

\paragraph{Nguyên nhân gốc rễ.}
v0: caller tự cấp \cmd{userId}/\cmd{cartId}/\cmd{itemId} mà không có danh tính xác
thực và không kiểm ownership, cho phép đọc/sửa/xóa giỏ của người khác (IDOR). Đây
là \cwe{639} (Authorization Bypass Through User-Controlled Key).

\paragraph{Bằng chứng lỗi tồn tại.}
Finding \cmd{VULN-005} với \cmd{integrity_impact = HIGH}. Đoạn v0 minh họa không
kiểm ownership:

\begin{lstlisting}[language=CSharp,numbers=none,breaklines=true,caption={v0 -- minh họa CartController thiếu ownership (theo VULN-005)}]
[HttpPost("user/{userId}/add")]
public async Task<IActionResult> AddToCart(int userId, [FromBody] AddToCartRequest request)
    => Ok(await _cartService.AddItemToCartAsync(userId, request.ProductId, request.Quantity)); // v0: khong kiem owner
\end{lstlisting}

\paragraph{Cách vá.}
JWT identity (claim \cmd{sub}) + \cmd{User.CanAccessUser(...)} trên mọi thao tác cart,
và \cmd{[Authorize]} ở cấp lớp.

\begin{lstlisting}[language=CSharp,numbers=none,breaklines=true,caption={v1 -- \texttt{CartController.cs}, ownership mọi thao tác}]
[ApiController]
[Route("api/[controller]")]
[Authorize]                       // v1: bat buoc xac thuc
public class CartController : ControllerBase
...
if (!User.CanAccessUser(userId)) return Forbid();          // v1: kiem so huu
...
if (!User.CanAccessUser(cart.UserId)) return Forbid();     // v1: kiem so huu
\end{lstlisting}

\paragraph{Bằng chứng đã vá.}
Integration test \cmd{Authentication_and_resource_ownership_are_enforced} kiểm
anonymous $\to$ 401, cross-user $\to$ 403 cho \texttt{api/cart/user/\{id\}}; HTTP
fuzz có 46 response 403. Giới hạn: chưa có test bắt buộc cho route cart mới.

% =====================================================================
%  PHẦN 5 -- BẢNG TỔNG HỢP BEFORE/AFTER
% =====================================================================
\clearpage
\section{Bảng tổng hợp before/after}

\footnotesize
\renewcommand{\arraystretch}{1.25}
\begin{longtable}{@{}L{1.3cm} L{3.2cm} L{3.5cm} L{3.5cm} L{2.4cm}@{}}
\caption{Bảng before/after bắt buộc. Mỗi dòng tương ứng đúng một mục ở Phần 4;
``v0/v1'' của C được chọn bằng \texttt{-UFIXED}/\texttt{-DFIXED}.}
\label{tab:beforeafter}\\
\toprule
\textbf{CWE} & \textbf{Vị trí} & \textbf{Trạng thái v0 (BUGGY)} &
\textbf{Trạng thái v1 (FIXED)} & \textbf{Công cụ xác minh} \\
\midrule
\endfirsthead
\multicolumn{5}{@{}l}{\small\itshape (tiếp theo Bảng~\ref{tab:beforeafter})}\\
\toprule
\textbf{CWE} & \textbf{Vị trí} & \textbf{Trạng thái v0 (BUGGY)} &
\textbf{Trạng thái v1 (FIXED)} & \textbf{Công cụ xác minh} \\
\midrule
\endhead
\bottomrule
\endlastfoot
CWE-401 & \path{order.c:6--18} &
LeakSanitizer: \texttt{detected memory leaks}, exit 1 &
Sạch leak, exit 0 &
ASan/LSan; CBMC; Valgrind \\
\addlinespace
CWE-415 & \path{payment.c:5--14} &
ASan double-free, exit $\neq$ 0 &
Exit 0, sạch &
ASan; CBMC; Valgrind \\
\addlinespace
CWE-416 & \path{cart.c:5--19} &
ASan UAF, exit $\neq$ 0 &
Exit 0; log đúng id &
ASan; CBMC; Valgrind \\
\addlinespace
CWE-122 & \path{product.c:6--20} &
ASan heap overflow &
Reject input dài; 5 ca biên sạch &
ASan; CBMC; Valgrind \\
\addlinespace
CWE-121 & \path{user.c:5--15} &
ASan stack overflow; corruption &
Reject input dài; 6 ca biên sạch &
ASan; CBMC; Valgrind \\
\addlinespace
CWE-190 & \path{PaymentVerification.z3} &
TEST6 \texttt{sat}: overflow int32 &
TEST7 \texttt{unsat}; chặn tổng/cent &
Z3 (7/7 ca) \\
\addlinespace
CWE-841 & \texttt{model BUGGY (JSON)} &
P1 false: PAID\_AFTER\_CANCELLED &
15/15 property; guard + transition + idempotency &
Model; TLC; 35 integration \\
\addlinespace
CWE-306 & \path{OrdersController.cs} &
Không có \cmd{[Authorize]}; đọc/hủy đơn bất kỳ ai &
JWT + \cmd{[Authorize]} + ownership &
Integration; HTTP (401/403) \\
\addlinespace
CWE-602 & \path{OrderService.cs} &
Lưu \cmd{UnitPrice}/\cmd{TotalAmount} từ client &
Giá lấy từ server + tự tính tổng + transaction &
Unit; SMT; HTTP \\
\addlinespace
CWE-20 & \path{InputValidator.cs} &
Lệch tổng $> 0{,}01$ rơi vào khối rỗng &
Từ chối mismatch và giá dưới ngưỡng &
Unit; boundary; fuzz \\
\addlinespace
CWE-639 & \path{CartController.cs} &
\cmd{userId}/\cmd{cartId} do client cấp &
JWT identity + ownership mọi thao tác &
Integration; HTTP (403) \\
\end{longtable}

% =====================================================================
%  PHẦN 6 -- BẰNG CHỨNG KIỂM CHỨNG
% =====================================================================
\clearpage
\section{Bằng chứng kiểm chứng}

\subsection{Tổng hợp các lần chạy}
\begin{table}[H]
\centering
\footnotesize
\renewcommand{\arraystretch}{1.2}
\begin{tabularx}{\textwidth}{@{}L{3.8cm} Y L{5.0cm}@{}}
\toprule
\textbf{Hạng mục} & \textbf{Kết quả} & \textbf{Evidence} \\
\midrule
Unit/regression & 149/149 PASS, 0 fail, 0 skip & \path{results/application_tests.results.json} \\
Integration & 35/35 PASS, 0 fail, 0 skip & \path{results/application_tests.results.json} \\
Recovery model & 14 state, 35 transition, 7 contract, 15 property, 6/6 mutation & \path{results/order_payment_recovery_model.results.json} \\
Đặc tả workbook & 10 sheet, 27 yêu cầu, 18 state, 10 CTL, 50 testcase & \path{results/teacher_spec_verification.json} \\
TLC & 35 state sinh, 20 phân biệt, depth 9 & \path{results/tlc/OrderPayment.tlc.results.json} \\
Payment SMT & 7/7 đúng kỳ vọng & \path{results/payment_verification.results.json} \\
Inventory SMT & 3/3, witness stock=10, q1=q2=6 & \path{results/inventory_verification.results.json} \\
C dynamic & 22/22 PASS & \path{c-memory/evidence/dynamic-analysis.json} \\
CBMC / Valgrind & 10/10 và 10/10 & \path{c-memory/evidence/bmc-analysis.json}, \path{c-memory/evidence/valgrind-analysis.json} \\
SAT/DRAT & 10/10, UNSAT core kiểm lại & \path{c-memory/evidence/sat-pipeline.json} \\
Fuzz .NET & 4/4, line 91,05\%, branch 97,50\% & \path{results/application_fuzzing.results.json} \\
Validator coverage & line 85,71\%, branch 97,22\% & \path{results/coverage_summary.json} \\
HTTP fuzz & 601 request, 108 feature, 500 semantic edge, 0 lỗi 5xx & \path{results/http_fuzz.results.json} \\
Browser CDP & 8 scenario, 7/7 fault hiện error state & \path{results/browser_network_fuzz.results.json} \\
C coverage fuzz & 20.000 mutation, edge 81$\rightarrow$148, queue 46, 0 crash & \path{results/c_coverage_fuzz.results.json} \\
Recovery & 2/2, ready 200$\rightarrow$503$\rightarrow$200 & \path{results/postgres_reconnect_*.results.json} \\
\bottomrule
\end{tabularx}
\caption{Tổng hợp bằng chứng. Đường dẫn tính từ \texttt{security-analysis/}.}
\end{table}

\subsection{Bằng chứng an toàn bộ nhớ}
Bảng~\ref{tab:cppcheck} đối chiếu số lượng cảnh báo trước/sau. Cppcheck chỉ bắt
được lỗi biên của \texttt{user.c} và một cảnh báo style; các lỗi còn lại do pattern
scanner chuyên biệt phát hiện. Đây là lý do nhóm dùng nhiều công cụ thay vì tin
tuyệt đối vào một công cụ.

\begin{table}[H]
\centering
\footnotesize
\renewcommand{\arraystretch}{1.2}
\begin{tabularx}{\textwidth}{@{}Y L{4.2cm} L{4.2cm}@{}}
\toprule
\textbf{Chỉ số} & \textbf{v0 (BUGGY)} & \textbf{v1 (FIXED)} \\
\midrule
Pattern scanner (5 rule) & 5 finding (1/module) & 0 finding \\
Cppcheck error & 1 (\texttt{bufferAccessOutOfBounds}, user.c:13) & 0 \\
Cppcheck style & 1 (\texttt{unusedStructMember}) & 1 (không phải lỗi bảo mật) \\
GCC \cmd{-Wall -Wextra -Wpedantic} & 0 diagnostic & 0 diagnostic \\
ASan/UBSan & 10/10 ca đúng diagnostic & 10/10 ca sạch \\
CBMC & 5 \texttt{VERIFICATION FAILED} đúng lỗi & 5 \texttt{VERIFICATION SUCCESSFUL} \\
Valgrind & 5 lỗi đúng loại & 5 sạch \\
\bottomrule
\end{tabularx}
\caption{Đối chiếu cảnh báo tĩnh và động trước/sau khi vá.}
\label{tab:cppcheck}
\end{table}

Pipeline không dùng clang-tidy; bộ phân tích tĩnh gồm Cppcheck, compiler
\cmd{-Wall -Wextra -Wpedantic} và pattern scanner viết riêng cho 5 lớp lỗi.
Cppcheck chỉ bắt 1/5 lỗi bộ nhớ, nên pattern scanner và sanitizer là bắt buộc
để đủ bằng chứng.

\begin{table}[H]
\centering
\footnotesize
\renewcommand{\arraystretch}{1.15}
\begin{tabularx}{\textwidth}{@{}L{2.6cm} L{1.0cm} L{6.2cm} Y@{}}
\toprule
\textbf{Case} & \textbf{Exit} & \textbf{Diagnostic ghi nhận} & \textbf{Kết luận} \\
\midrule
\cmd{order_buggy} & 1 & \texttt{LeakSanitizer: detected memory leaks} & CWE-401 tái hiện \\
\cmd{order_fixed} & 0 & clean sanitizer run & CWE-401 đã vá \\
\cmd{payment_buggy} & 1 & \texttt{AddressSanitizer: attempting double-free} & CWE-415 tái hiện \\
\cmd{payment_fixed} & 0 & clean sanitizer run & CWE-415 đã vá \\
\cmd{cart_buggy} & 1 & \texttt{AddressSanitizer: heap-use-after-free} & CWE-416 tái hiện \\
\cmd{cart_fixed} & 0 & clean sanitizer run & CWE-416 đã vá \\
\cmd{product_buggy} & 1 & \texttt{AddressSanitizer: heap-buffer-overflow} & CWE-122 tái hiện \\
\cmd{product_fixed} & 0 & clean sanitizer run & CWE-122 đã vá \\
\cmd{user_buggy} & 1 & \texttt{AddressSanitizer: stack-buffer-overflow} & CWE-121 tái hiện \\
\cmd{user_fixed} & 0 & clean sanitizer run & CWE-121 đã vá \\
\bottomrule
\end{tabularx}
\caption{10 ca lõi trong bộ 22 ca sanitizer. 12 ca còn lại là các độ dài biên của
product/user, tất cả exit 0 và sạch sanitizer.}
\label{tab:sanitizer}
\end{table}

\anh{anh-01-sanitizer.png}{Bộ 22 ca sanitizer PASS; mỗi ca BUGGY có đúng
diagnostic, mỗi ca FIXED sạch.}

\anh{anh-02-cwe401-truoc-sau.png}{Trước/sau CWE-401: bản BUGGY rò rỉ, bản FIXED sạch.}

\anh{anh-03-cbmc-valgrind.png}{CBMC
(bound 64) và Valgrind Memcheck đều 10/10.}

\subsection{Bằng chứng mô hình, SMT và TLC}
\anh{anh-04-smt-tlc.png}{SMT payment 7/7 và TLC
35 trạng thái, depth 9.}

\anh{anh-10-model-witness.png}{Mô hình v0 để lộ chuyển
tiếp bất hợp lệ (đối chiếu Hình~\ref{fig:order-state}); bản v1 chặn tại chỗ.}

\anh{anh-05-sat-drat.png}{Pipeline
CNF/DIMACS + CaDiCaL + DRAT + UNSAT core đạt 10/10.}

\subsection{Bằng chứng fuzzing và coverage}
\anh{anh-06-fuzz-coverage.png}{Bộ 4 test fuzz .NET 4/4 và coverage
dòng/nhánh của 5 target.}

\anh{anh-07-http-fuzz.png}{Fuzz HTTP black-box trên API + PostgreSQL thật.}

\anh{anh-09-browser-cdp.png}{Chrome thật bị chặn response \texttt{/api/products}; 7/7 fault hiện error state + retry.}

\subsection{Bằng chứng phục hồi}
\anh{anh-08-recovery.png}{Charge/refund phục hồi sau PostgreSQL stop/start,
cùng process API, không gọi lại gateway.}

\subsection{Bảng seed và số lần chạy}
\begin{table}[H]
\centering
\small
\renewcommand{\arraystretch}{1.2}
\begin{tabularx}{\textwidth}{@{}L{4.3cm} L{2.2cm} L{3.4cm} Y@{}}
\toprule
\textbf{Campaign} & \textbf{Seed} & \textbf{Budget} & \textbf{Kết quả} \\
\midrule
.NET validator & 20261009 & 25.000 $\times$ 2 validator & 0 crash, 0 vi phạm invariant \\
.NET stateful & 20261009 & 2.000 chuỗi $\times$ 20 bước & Invariant Order/Payment/Inventory giữ \\
HTTP fuzz & 20261022 & 500 mutation (601 request) & 0 lỗi 5xx, 0 vi phạm DB \\
C coverage-guided & 20261022 & 20.000 mutation & 0 crash, edge 81$\rightarrow$148 \\
Browser CDP & -- & 8 scenario & 7/7 fault hiện error state \\
\bottomrule
\end{tabularx}
\caption{Các campaign và seed. Seed cố định để mọi kết quả tái lập được.}
\end{table}

Bảng trên là các lần chạy đã lưu trong evidence của nhóm. Dự án chọn seed cố định
để tái lập chính xác; nếu cần bổ sung nhiều seed, lặp lại lệnh với
\texttt{-{}-seed} khác và ghi kết quả vào cùng bảng, không thay đổi oracle.

% =====================================================================
%  PHẦN 7 -- PHẢN BIỆN & BÀI HỌC
% =====================================================================
\section{Phản biện \& Bài học}

\subsection{Nếu bắt đầu lại, nhóm sẽ thiết kế lại phần nào?}
Nhóm sẽ đưa state machine và idempotency vào ngay từ đầu thay vì bổ sung sau:
mọi thao tác tạo payment phải đi qua một bảng chuyển tiếp duy nhất, khóa
idempotency phải là duy nhất ở tầng cơ sở dữ liệu, và tiền phải là \cmd{decimal}
với ràng buộc cent ngay ở DTO. Nhóm cũng sẽ viết kịch bản crash-injection sớm hơn
để buộc thiết kế recovery phải rõ ràng trước khi tối ưu hiệu năng.

\subsection{Có lỗ hổng nào nhóm nghi ngờ vẫn còn nhưng chưa xác nhận?}
Có. Các cửa sổ crash chưa được mô phỏng: mất điện đúng lúc ghi journal, partition
giữa API và PostgreSQL, webhook đến trễ hoặc trùng từ provider thật. Rate limiter
và audit retention hiện chạy trong bộ nhớ/CSDL cục bộ, chưa kiểm với nhiều instance.
Các lỗi C trong báo cáo là mô phỏng didactic; chúng không chứng minh runtime C\#
có cùng lỗi, và bản thân mô hình SMT/TLC cũng có bound.

\subsection{Bài học lớn nhất về phương pháp luận}
Bài học lớn nhất là \emph{evidence-first}: mỗi khẳng định phải gắn với một artifact
có thể chạy lại (seed, log, hash), và mỗi checker phải được kiểm bằng mutation âm.
Nhóm học được cách phân biệt ``SAFE up to bound'' với chứng minh vô hạn, và cách
dùng nhiều công cụ bổ trợ vì không công cụ nào bắt hết lỗi: Cppcheck bỏ sót 4/5 lỗi
C, trong khi ASan chỉ thấy lỗi khi chương trình thực sự chạy đúng nhánh.

% =====================================================================
%  PHẦN 8 -- PHỤ LỤC
% =====================================================================
\clearpage
\section{Phụ lục}

\subsection*{A. Hướng dẫn build và chạy}
Yêu cầu: .NET 10 SDK, Python 3.10+, GCC/Make, Cppcheck, CBMC, Valgrind, Z3,
Java + \texttt{tla2tools.jar}, Chrome cho test browser; PostgreSQL binaries cho
integration/recovery; CaDiCaL + drat-trim cho pipeline SAT. Chạy từ gốc repository.

Cài TLA+ một lần (nếu chưa có \texttt{tmp/tools/tla2tools.jar}):
\begin{lstlisting}[language=bash,caption={Tải tla2tools.jar cho TLC}]
mkdir -p tmp/tools
curl --fail --location --retry 2 \
  https://github.com/tlaplus/tlaplus/releases/download/v1.7.4/tla2tools.jar \
  --output tmp/tools/tla2tools.jar
\end{lstlisting}

\begin{lstlisting}[language=bash,caption={Build và chạy kiểm chứng từ gốc repository}]
dotnet restore OrderPaymentMicroservice.slnx
dotnet build OrderPaymentMicroservice.slnx --no-restore --nologo

dotnet test OrderPaymentMicroservice.Tests/OrderPaymentMicroservice.Tests.csproj \
  --no-restore --nologo
dotnet test OrderPaymentMicroservice.IntegrationTests/OrderPaymentMicroservice.IntegrationTests.csproj \
  --no-restore --nologo --blame-hang-timeout 90s

python3 security-analysis/verify_recovery_model.py
python3 security-analysis/audit_liveness.py
python3 security-analysis/verify_threat_model.py
python3 security-analysis/verify_teacher_spec.py

make -C security-analysis/c-memory test
make -C security-analysis/c-memory static-check
make -C security-analysis/c-memory bmc-check
make -C security-analysis/c-memory valgrind-check
make -C security-analysis/c-memory sat-check

python3 security-analysis/verify_payment_smt.py
python3 security-analysis/verify_inventory_smt.py
python3 security-analysis/run_tlc.py --jar tmp/tools/tla2tools.jar --timeout 120

python3 security-analysis/run_fuzzing.py
python3 security-analysis/http_fuzz.py --iterations 500 --no-build
python3 security-analysis/verify_checkout.py
python3 security-analysis/browser/network_fuzz.py
python3 security-analysis/run_c_coverage_fuzz.py --iterations 20000
python3 security-analysis/run_recovery_demo.py

dotnet test OrderPaymentMicroservice.slnx --no-restore --nologo \
  --logger 'trx;LogFilePrefix=verification' \
  --results-directory security-analysis/results/dotnet \
  --blame-hang-timeout 90s
\end{lstlisting}

\subsection*{B. Danh mục evidence}
\begin{table}[H]
\centering
\footnotesize
\renewcommand{\arraystretch}{1.2}
\begin{tabularx}{\textwidth}{@{}L{5.0cm} Y@{}}
\toprule
\textbf{Tệp} & \textbf{Nội dung} \\
\midrule
\path{results/application_tests.results.json} & 149 unit + 35 integration, 0 fail/skip \\
\path{results/order_payment_recovery_model.results.json} & 14 state, 35 transition, 7 contract, 15 property, 6 mutation \\
\path{results/order_payment_state_response_matrix.csv} & Ma trận state $\rightarrow$ response \\
\path{results/teacher_spec_verification.json} & 10 sheet, 27 yêu cầu, 50 testcase, CWE-667/693 cố ý \\
\path{results/threat_model_results.json} & 6 asset, 4 boundary, 9 threat, 6/6 STRIDE \\
\path{results/liveness_audit.results.json} & raw AG AF false, progress-only true \\
\path{results/payment_verification.results.json} & 7/7 SMT, TEST6 sat, TEST7 unsat \\
\path{results/inventory_verification.results.json} & 3/3 SMT, witness stock=10, q1=q2=6 \\
\path{c-memory/evidence/dynamic-analysis.json} & 22 ca sanitizer \\
\path{c-memory/evidence/static-analysis.json} & pattern scanner, cppcheck, layout \\
\path{c-memory/evidence/bmc-analysis.json} & 10 ca CBMC, bound 64 \\
\path{c-memory/evidence/valgrind-analysis.json} & 10 ca Memcheck \\
\path{c-memory/evidence/sat-pipeline.json} & 10 ca CNF/DRAT/UNSAT core \\
\path{results/application_fuzzing.results.json} & 4/4 fuzz, coverage 91,05/97,50 \\
\path{results/http_fuzz.results.json} & 601 request, 0 lỗi 5xx \\
\path{results/browser_network_fuzz.results.json} & 8 scenario, 7/7 fault UI \\
\path{results/postgres_reconnect_*.results.json} & charge/refund 200$\rightarrow$503$\rightarrow$200 \\
\path{results/process_restart.results.json} & kill API, phục hồi charge \\
\path{results/refund_restart.results.json} & kill API, phục hồi refund \\
\path{results/evidence_manifest.json} & Hash source/evidence, commit, dotnet version \\
\bottomrule
\end{tabularx}
\caption{Danh mục evidence chính; đường dẫn tính từ \texttt{security-analysis/}. Log đầy đủ nằm trong repository, không cắt bớt.}
\end{table}

\subsection*{C. Phân công công việc}
\begin{table}[H]
\centering
\small
\renewcommand{\arraystretch}{1.25}
\begin{tabularx}{\textwidth}{@{}L{4.4cm} L{3.2cm} Y@{}}
\toprule
\textbf{Thành viên} & \textbf{MSSV} & \textbf{Công việc chính} \\
\midrule
Phạm Văn Hoàng & 23010245 & Toàn bộ: đặc tả, threat model, API Order/Payment, 5 module C, SMT/TLC, fuzz, recovery, báo cáo \\
\bottomrule
\end{tabularx}
\caption{Bảng phân công công việc.}
\end{table}

\subsection*{D. Mẫu commit message theo từng CWE}
Mỗi message nêu nguyên nhân, cách vá và bằng chứng đã chạy lại. Đây là mẫu tham
khảo bám theo thay đổi v1 trong repo, không phải log commit thật.

\subsubsection*{D.1. Năm lỗi bộ nhớ trong module C}
\begin{lstlisting}[language=bash,caption={Commit mẫu cho CWE-401/415/416/122/121.}]
fix(CWE-401): free owned items before parent in cancel_order
- Free order->items before free(order); no ownership transfer
- Add LeakSanitizer regression case order_fixed
Verified: ASan/LSan clean (0 leaks; was 240 bytes in 5 allocations)

fix(CWE-415): null transaction_code after first free
- Set pointer to NULL so common cleanup free(NULL) is a no-op
Verified: ASan clean (was: attempting double-free)

fix(CWE-416): read product_id before freeing removed cart item
- Copy product_id, then free(removed) and set removed = NULL
- Log from the copy instead of the freed pointer
Verified: ASan clean (was: heap-use-after-free)

fix(CWE-122): reject product name of 32 bytes or more
- Check strlen(input) >= NAME_CAPACITY before strcpy
- Add boundary cases 0/1/31/32/128
Verified: ASan clean (was: heap-buffer-overflow)

fix(CWE-121): reject oversized email before strcpy
- Check strlen(input) >= sizeof(email); reject instead of truncate
- Add boundary cases 0/1/31/32/128
Verified: ASan/CBMC/Valgrind clean (was: stack-buffer-overflow)
\end{lstlisting}

\subsubsection*{D.2. Lỗi thiết kế và nghiệp vụ trong API .NET}
\begin{lstlisting}[language=bash,caption={Commit mẫu cho CWE-190/841/306/602/20/639.}]
fix(CWE-190): validate cents and use checked arithmetic for totals
- Require two-decimal money values via MoneyAmounts.RequireCents
- Compute subtotal and total with checked(); reject overflow
Verified: SMT payment 7/7 PASS; arithmetic unit tests PASS

fix(CWE-841): enforce order state whitelist in Transition
- Replace free-form status writes with CanTransition(from, to)
- Reject CANCELLED -> PAID and every other illegal edge
Verified: TLC 35 states/depth 9; state-machine tests PASS

fix(CWE-306): require authentication on order and cart endpoints
- Add [Authorize]; admin-only for confirm, complete and list-all
- Derive customer id from User.GetRequiredUserId(), not from body
Verified: 401/403 integration test PASS; HTTP fuzz: 601 requests, no 5xx

fix(CWE-602): compute order total server-side from product prices
- Ignore client TotalAmount; reload Product.Price and recompute
- Reject mismatch above 0.01; never trust client unit price
Verified: ClientCannotCreateOrderUsingTamperedPrice PASS

fix(CWE-20): validate order bounds, quantities and totals
- Bound item count, quantity 1..10000, price/total 0.01..999999.99
- Check item subtotal and total consistency before persistence
Verified: MismatchOrBelowMinimumPriceIsRejected PASS; boundary cases PASS

fix(CWE-639): enforce ownership before cart/order access
- Call User.CanAccessUser(ownerId) on every by-id read and write
- Return 403 instead of exposing another user's resource
Verified: Authentication_and_resource_ownership_are_enforced PASS
\end{lstlisting}

\subsection*{E. Quy ước phiên bản và ánh xạ}
Repo không tách \texttt{*\_v0.c}/\texttt{*\_v1.c} mà dùng tiền xử lý:
\cmd{-UFIXED} cho v0 (BUGGY) và \cmd{-DFIXED} cho v1 (FIXED). Hai model
\texttt{order\_state\_model.json} và \texttt{order\_state\_model\_buggy.json} tương
ứng v1/v0 ở tầng thiết kế. Khi nộp kèm mã nguồn, cần giữ nguyên cả hai biến thể và
ghi rõ dòng nào thuộc biến thể nào; không sửa file gốc sau khi đã sinh evidence.

% =====================================================================
%  TÀI LIỆU THAM KHẢO
% =====================================================================
\clearpage
\section*{Tài liệu tham khảo}
\addcontentsline{toc}{section}{Tài liệu tham khảo}
\begin{enumerate}[leftmargin=1.6em]
\item Hướng dẫn Viết Báo cáo Đồ án, CSE703153 -- Kỹ thuật Lập trình An toàn,
Học viện Công nghệ Thông tin, Đại học Phenikaa, 2026.
\item E. Clarke, D. Kroening, and F. Lerda, ``A Tool for Checking ANSI-C
Programs,'' TACAS, 2004.
\item L. de Moura and N. Bjørner, ``Z3: An Efficient SMT Solver,'' TACAS, 2008.
\item L. Lamport, \textit{Specifying Systems: The TLA+ Language and Tools for
Hardware and Software Engineers}, Addison-Wesley, 2002.
\item D. Kroening and M. Tautschnig, ``CBMC -- Bounded Model Checking for
Software,'' TACAS, 2014.
\item MITRE, \textit{CWE -- Common Weakness Enumeration}, \url{https://cwe.mitre.org/}.
\item SEI CERT C Coding Standard, \url{https://wiki.sei.cmu.edu/confluence/display/c/}.
\item Mã nguồn và evidence của nhóm:
\mustfill{điền URL repository nếu nộp kèm}.
\end{enumerate}

\end{document}
